\chapter{TEST LOG}
\label{ch:tests}

This chapter records the tests run on the robot, what was measured and what
was concluded. The largest part is the investigation of the board resets
that blocked the chase tests, which ended with the camera on its own power
through a USB-C hub (\S\ref{sec:resettests}).

\section{SUMMARY OF EARLIER TESTS}

Table~\ref{tab:earliertests} lists the bring-up measurements that are
described in their own chapters.

\begin{table}[htbp]
\centering
\caption{Bring-up tests described elsewhere}
\label{tab:earliertests}
\small
\begin{tabularx}{\linewidth}{>{\raggedright\arraybackslash}p{38mm}>{\raggedright\arraybackslash}X>{\raggedright\arraybackslash}p{18mm}}
\toprule
Test & Result & Section \\
\midrule
Encoder push, \SI{0.914}{m} (3 ft) & \SI{0.914}{m} measured after setting 15 counts per wheel turn & \S\ref{sec:odom} \\
Motor map, wheels lifted & Motor numbering and signs found; about \SI{0.4}{m/s} at 30\% duty & \S\ref{sec:duty} \\
Gyro drift, parked & \SI{1.6}{\degree/min} before the zero-velocity update, zero after & \S\ref{sec:zupt} \\
Magnetometer, two-turn spin & Circle fit offsets; compass agrees with the gyro to 3\% & \S\ref{sec:mag} \\
Lidar mount test & Zero bearing points to the rear; transform yaw set to \SI{180}{\degree} & \S\ref{sec:slamdiag} \\
Turn rate on the floor & \SI{5.0}{rad/s} commanded gave \SI{1.31}{rad/s}; effective track \SI{0.46}{m} & \S\ref{sec:duty} \\
Camera calibration & Epipolar error 0.36~px, baseline \SI{7.43}{cm} & \S\ref{sec:oakcal} \\
\bottomrule
\end{tabularx}
\end{table}

\section{CHASE, DRY RUN AND WHEELS OFF THE GROUND}
\label{sec:chasetest}

With the merged chase fix (\S\ref{sec:chaseyaw}) on the board and the camera
on a USB Y-cable fed by a battery bank, chase was run with
\texttt{targets:=[dog, person]} on 3 October 2026.

\paragraph{Dry run.} Nothing was published to the motors. The node tracked
both the dog (64--98\% confidence) and a person (61--99\%) at
\SIrange{0.5}{2.4}{m}. The forward command fell to zero inside the
\SI{0.9}{m} follow distance and reached about \SI{0.19}{m/s} at
\SI{1.4}{m}, as \eqref{eq:chasev} predicts. Stepping to the robot's right
gave a negative turn command, the correct sign, and bearings inside the
\SI{5}{\degree} dead band gave none.

\paragraph{Driving, wheels lifted.} With \texttt{dry\_run:=false} and R2 held,
the node logged \texttt{Detector sees 39.2 deg horizontally (fx=1516 px at
1080p, centre crop)}, confirming the preview crop of
\S\ref{sec:previewcrop}. Turn requests of about \SI{0.3}{rad/s} were raised
by the yaw-rate loop to between 0.7 and \SI{3.8}{rad/s}. This is the
expected behaviour with the wheels in the air: the gyro sees no rotation,
so the integrator keeps adding command until the limit. Releasing R2 returned the node to
\texttt{idle}. A few minutes into this run the
board reset, which started the investigation below.

\section{BOARD RESET INVESTIGATION}
\label{sec:resettests}

\subsection{What had been ruled out}

After the reset during the chase drive, the previous boot's journal and the
serial console showed no kernel panic, no thermal shutdown message and no
watchdog activity (the watchdogs were off). The CPU had been well below its
\SI{85}{\degree C} critical trip point. So the reset was a hardware power
event, not a software crash. The once-per-second battery logger used until
then was too slow to show a transient, and it did not restart after a
reboot.

\subsection{Fast logger}
\label{sec:fastlog}

A new logger (\file{\textasciitilde/bin/fastlog.py} on the board) was written for
the tests:
\begin{itemize}
\item It reads the battery pack and DC jack voltages from the cape's ADC in
  a tight loop at about \SI{250}{Hz}, and writes the minimum, mean and
  maximum of each channel for every \SI{50}{ms} window, so a dip shorter than
  a second still shows.
\item Every \SI{0.5}{s} it also logs every thermal zone (CPU, GPU, cores),
  the CPU load and the CPU clock.
\item A second process follows \texttt{dmesg -w} into its own file.
\item Every line is flushed and \texttt{fsync}ed as it is written, so the
  record survives a hard reset up to the last few milliseconds.
\end{itemize}
The \SI{5}{V} and \SI{3.3}{V} rails are not wired to any ADC channel on this
board, so they could not be logged; a multimeter on SYS\_5V (P9.7/P9.8) is
the only way to see them.

\subsection{Isolation tests}

Each test isolated one load. The results are in
Table~\ref{tab:resettests}.
\begin{description}
\item[Test A, motors only.] With the wheels off the ground and the camera
  not running, \topic{/cmd\_vel} was published at \SI{10}{Hz} in a repeating
  pattern: \SI{15}{s} forward at full speed (the base caps the duty at 0.6),
  \SI{15}{s} reverse, \SI{10}{s} spins at $\pm\SI{10}{rad/s}$, and abrupt
  forward/reverse flips every second.
\item[Test B, camera only, robot running.] \texttt{robot.launch} up (lidar,
  base, IMU), no motor commands, and the full chase pipeline (stereo depth
  plus detector) running in dry run.
\item[Test B2, camera only, everything else off.] The robot launch file was
  stopped, so no lidar, no motor driver and no ROS ran; only the camera
  pipeline.
\end{description}

\begin{table}[htbp]
\centering
\caption{Board reset isolation tests (3 October 2026)}
\label{tab:resettests}
\small
\begin{tabularx}{\linewidth}{lllX}
\toprule
Test & Camera & Motors and lidar & Result \\
\midrule
A & off & full duty, \SI{6.5}{min} & No reset. \\
B & streaming & lidar on, motors idle & Reset \SI{17}{s} into camera streaming. \\
B2 & streaming & all off & Reset at \SI{2}{min} \SI{13}{s}. \\
Final & on USB-C hub with battery bank & all on & Stable. \\
\bottomrule
\end{tabularx}
\end{table}

\paragraph{Observations.}
\begin{itemize}
\item The pack and DC jack voltages were flat right up to the last
  \SI{50}{ms} window before each reset: no dip on the cape's input.
\item There were no USB errors and no kernel messages before the resets.
\item The SoC ran at \SIrange{62}{70}{\degree C}, with the CPU throttled to
  \SI{400}{MHz} throughout the camera runs. Hot, but a thermal shutdown would
  have left a message.
\item Whenever the motors ran, the battery pack reading fell to about
  \SI{1.3}{V} and recovered when they stopped, while the DC jack held
  \SI{8.8}{V}. That looks like the pack's protection circuit cutting out
  under motor current while the supply on the DC jack carried the load.
  It did not cause a reset, but on the battery alone the motors may not run
  at full duty. \emph{Unconfirmed.}
\end{itemize}

\subsection{Conclusion}

The camera is the trigger and the motors are not: the motors at full duty
for six and a half minutes caused nothing, and the camera alone reset the
board with nothing else running. Since the cape's input stayed flat, the
failure is downstream of it, on the \SI{5}{V} path that feeds the
BeagleBone and its USB-A port. That agrees with the power budget
(Table~\ref{tab:budget}): the board and the camera alone draw about
\SIrange{1.6}{2.7}{A} from a \SI{2}{A} regulator, so the full load was never
needed to exceed it. The earlier \SI{90}{s} camera-only run that passed
(\S\ref{sec:power}) was simply too short; test B2 took over two minutes.

The USB Y-cable did not fix it. Most Y-cables still connect the host's
\SI{5}{V} line to the device, so the camera could keep drawing from the board
whenever the bank's output was lower or switched off.

\subsection{The fix}

Jason moved the camera to a USB-C hub plugged into the BeagleBone's USB-C
port, with the battery bank plugged into the hub's power input. The camera's
current now comes from the bank through the hub instead of from the cape's
\SI{5}{V} rail. With this arrangement the whole system, including the
camera, has run stably. The USB-C port had not powered the camera in the
first tests; through the powered hub it does.

\influx{The current drawn from the bank, and whether the hub also feeds the
BeagleBone through its USB-C port, have not been measured. A long run with
chase driving on the floor, and a run on the robot's battery alone, are still to be logged.}
